
\textit{Fully Homomorphic Encryption (FHE)}~\cite{gentry2009fully, brakerski2014leveled, halevi2019BFVimproved, kim2022approximate} is a cryptographic privacy-preserving computation technique that enables a unique computation model where mutually-distrusting parties can collaborate. For instance, computation providers, such as Machine Learning as a Service (MLaaS) model operators or financial analysis model proprietors, aim to offer services to clients without revealing their proprietary models. Simultaneously, clients seek to supply data for computations while keeping the contents confidential. 
FHE facilitates this cooperative model by allowing the server to process the client's encrypted data directly without first decrypting it. This ensures the confidentiality of both the client's data and the output. Therefore, FHE acts as a dual-purpose safeguard, securing sensitive client data as well as proprietary models, making it a crucial tool for enabling privacy-protected computations in cloud environments~\cite{AmazonScience, IBMSecurity, GoogleCloud, Google, Microsoft, reagen2021cheetah}.

FHE has seen a surge of interest in the computing industry in recent years. FHE has been deployed in learning-enabled systems to build privacy-preserving machine learning inferences, deployed on real-world products including Duality Techniques~\cite{DualityTech}, Cape Privacy~\cite{capeprivacy}, Zama~\cite{Zama}, and Inpher~\cite{Inpher}. Some of FHE's weaknesses are being addressed. For example, its high execution time has improved substantially and over the past two years, the latency for deep neural network inference on ResNet-20 has decreased from 88,320 seconds to just 15.9 seconds~\cite{fan2023tensorfhe,samardzic2022craterlake, kim2022ark}, a phenomenal progress which will likely continue.


While FHE ensures data privacy, it ensures neither {\em computation integrity} (the property that code and data are not tampered with as they are used) nor {\em verifiability} (the property that the correct computation was performed). Verifiability includes computation integrity but adds the requirement that the code executed implements the correct computation. It is critical for a client to be able to verify the correctness of the calculation performed by the server. An adversary may compromise service to modify computation or user data,  leading to wrong output. Alternatively, an adversary may compromise the system (OS or hypervisor) or even indirectly modify data in memory without code compromise, e.g., via a RowHammer attack. This means that the client may unknowingly receive the wrong computation results. 

Recognizing this limitation, a cryptographic approach for integrity-checking protocols has been developed to address these integrity challenges and complement FHE. These include homomorphic message authentication codes (MAC, as discussed in \cite{chatel2022verifiable}) and zero-knowledge proofs (ZKP, as outlined in \cite{viand2023verifiable}). While these methods assure the integrity of computations on encrypted data, they do not ensure verifiability independently and introduce significant computational overhead due to their complexity regarding computational requirements and data volume. This leads to substantial processing delays and increased resource demands, posing a challenge for practical implementation ~\cite{bhadauria2020ligero++, bunz2018bulletproofs, gennaro2010non, goldwasser2015delegating, parno2016pinocchio, weng2021wolverine, brakerski2011fully, bois2021flexible, fiore2014efficiently, fiore2020boosting, ganesh2021rinocchio}.

The computation providers could also run their models in Trusted Execution Environments (TEEs). However, relying on TEEs is problematic for several reasons. First, TEEs do not inherently ensure verifiability, as the client cannot automatically verify if the correct computation was performed by the TEE~\footnote{An attestation can solve this but it requires the client to trust the compute platform correctness, the attester, and the correct thing is being attested.}  Second, TEE support is not universal; platforms such as NPUs and many GPUs do not support TEEs. Finally, some TEEs may not support full integrity protection even when TEE is supported. For example, AMD SEV~\cite{AMD_SEV} and Intel Scalable SGX~\cite{Intel_SGX} only protect the integrity of one enclave against other enclaves but not against direct memory tampering incurred by, say, RowHammer attacks. Only Intel SGX supports full integrity protection via an integrity tree~\cite{intel2021sgx}. 

Algorithm-Based Fault Tolerance (ABFT) is a method developed to verify computation results by embedding redundancy within data structures. Primarily used to detect and correct errors in high-performance computing, ABFT employs checksum techniques in matrix operations by adding extra rows or columns \cite{abft,abft1,abft2,ABFT-1984-checksum}. These checksums adapt during computations to identify discrepancies from faults. While effective for hardware failures or unintentional errors, ABFT is not designed to ensure data confidentiality or protect against intelligent adversaries that may manipulate data and checksum simultaneously to evade detection. 

In this paper, we propose \textit{DataSeal}, an approach that enables the \textit{verifiability} of FHE computations through the strategic integration of ABFT with low overheads. The core innovation of DataSeal lies in incorporating ABFT into the FHE framework by embedding secrets and redundancy within the original data. 
Similar to ABFT, the client can, at low cost, verify that the correct computation was performed by analyzing the redundancy in the resulting output produced by the cloud. Unlike ABFT, the redundant data is calculated using client-specific secret key, hence not having the key, the attacker cannot produce a fake checksum that evades verification. This synergy allows for the execution of FHE computations while allowing the verification of their computation results. The redundant data is also encrypted with FHE, ensuring confidentiality.  

To summarize, our contributions are:
\begin{enumerate}
\item A new approach (referred to as {\em DataSeal}), where we adapted ABFT for integration into FHE to enable computation verification while fortifying its security with a client-private key.  
\item The design of DataSeal for various matrix computations with low space overheads (i.e., two additional rows are added regardless of the matrix size), and scalable computation overheads (i.e. diminishing verification overheads as the matrix size increases). 
\item We demonstrate that DataSeal is secure against targeted data modifications. 
\end{enumerate}

We conduct performance evaluations for individual matrix operations (addition, multiplication, and convolution) and end-to-end DNN inference (including activation) to demonstrate the efficacy and practicality of our approach. The evaluation looks into the performance overheads on the server side. The evaluation of DataSeal shows its substantially lower overheads than other methods such as MAC, ZKP, and TEE, sometimes by one order of magnitude or higher. Furthermore, we found that DataSeal's verification overheads decrease relative to the computation as the problem size increases. In contrast, other schemes, such as FHE with MAC, FHE with ZKP, and FHE with TEE, their overheads tend to increase with the problem size. Thus, as the problem size increases, the performance gap between DataSeal and other schemes increases. For large enough matrix sizes, DataSeal overheads become nearly negligible.   This comparative analysis establishes DataSeal's unique attractiveness.

%We conduct evaluations for individual matrix operations (addition, multiplication, and convolution) and end-to-end DNN inference (including activation) to demonstrate the efficacy and practicality of our approach. The evaluation looks into the performance overheads on the server side. The evaluation shows that DataSeal maintains a constant overhead of two additional matrix rows. This results in an overhead of 2X for a single row vector but scales efficiently for larger matrix data, reducing the overhead of providing verification and data integrity in FHE by up to 99\% compared with the hardware Trusted Execution Environment (TEE) approach and software cryptographic approach. These results underscore the minimal computational burden introduced by our approach and highlight its superiority in client and server environments. This comparative analysis establishes our method's practicality and effectiveness in real-world applications, offering a more efficient alternative to existing state-of-the-art methods.
